% ==================================================================================================================================
% Notions d'analyse matricielle

\section{Notions d'analyse Matricielle}

\subsection{Rappels}

Soient $E$ et $F$ deux espaces vectoriels normés. $ \mathcal{L}(E,F)$ désigne l'espace vectoriel des applications linéaires 
de E vers F. Il est muni d'une norme appelée \textbf{norme opérationnelle} : 
    \[ \forall f \in \mathcal{L}(E,F), \quad \| f \| = \min \{ C \in \R_+ \; | \; \forall x \in E, \| f(x) \|_F \leqslant C \| x \|_E \} \]

\begin{definition}[Norme Subordonnée]
    Si $\|.\|$ est une norme sur $\K^n$ alors la norme subordonnée associée à $\|.\|$ sur $ \mathcal{M}_n(\K)$ est 
    la norme opérationnelle associée à $\|.\|$ sur $\K^n = E$ et $\|.\|$ sur $\K^n = F$. On la note de la même façon. 
    \[ \text{i.e} \quad \forall M \in \mathcal{M}_n(\K), \quad  \| M \| = \min \{ C \in \R_+ \; | \; \forall x \in \K^n, \| Mx \| \leqslant C \| x \| \} \]
\end{definition}

\begin{proposition}
    Soient $M,N \in \mathcal{M}_n(\K)$ alors :
        \[ \| M \times N \| \leqslant \|M\| \times \|N\| \]
    Ceci d'écoule de l'inégalité de Cauchy-Schwarz. 
\end{proposition}

\subsection{Intégrale d'une matrice}

\begin{remark}
    Soit $f : I \longrightarrow \mathcal{M}_n(\K)$ où I est un intervalle. Supposons que $f$ soit continue. 
    On peut alors définir $ \int_I f(t) \; dt \in \mathcal{M}_n(\K)$ composante par composante telle que :
        \[ \left( \int_I f(t) \; dt \right)_{i,j} = \int_I f(t)_{i,j} \; dt \] 
    Autrement dit, l'intégrale d'une matrice d'applications continues sur un intervalle est exactement la matrice 
    les intégrales des applications. 
\end{remark}

\begin{proposition}
    Soit $f : I = ]a,b[ \overset{ \mathcal{C}^0}{\longrightarrow} \mathcal{M}_n(\K)$. Alors :
        \[ \left\| \int_I f(t) \; dt \right\| \leqslant \int_{a}^{b} \left\| f(t) \right\| \; dt \] 
\end{proposition}

\begin{prop}[Théorème Fondamental de l'Analyse]
    Dans ce cas ci, le théorème fondamental de l'analyse est toujours vrai. 
    Autrement dit, soit $f : I = ]a,b[ \overset{ \mathcal{C}^0}{\longrightarrow} \mathcal{M}_n(\K)$. 
    Alors : 
        \[ \int_{a}^{b} f'(t) \; dt = f(b) - f(a) \] 
\end{prop}

\begin{remark}
    Soit $f \in \mathcal{C}^1(I, \mathcal{M}_n(\K))$. Alors la formule $(f^2)' = 2f'f$ est fausse 
    puisque $\mathcal{M}_n(\K)$ n'est pas un anneau commutatif. D'où : 
        \[ (fg)' = f'g + fg' \quad \text{où} \quad g \in \mathcal{C}^1(I, \mathcal{M}_n(\K)) \] 
\end{remark}

\subsection{Exponentielle d'une matrice}

\begin{remark}
    Pour rappel, on sait que $ \forall z \in \C, \; e^z = \sum_{k=0}^{\infty} \frac{z^k}{k!}$. 
\end{remark}

\begin{definition}[Exponentielle d'une matrice]
    Soit $A \in \mathcal{M}_n(\K)$ on a alors :
        \[ \exp(A) := \sum_{k=0}^{\infty} \frac{A^k}{k!} \]
    Cette série converge toujours. 
\end{definition}

\begin{proposition}[Propriétés de l'exponentielle]
    Soient $A,B \in \mathcal{M}_n(\K)$ alors 
        \[ e^{A+B} = e^A e^B \iff AB = BA \] 
\end{proposition}

\begin{proposition}[Inverse]
    Soit $A \in \mathcal{M}_n(\K)$ alors $e^A$ est inversible. 
    Pour cela, on pose :
    \[ e^0 = \sum_{k=0}^{\infty} \frac{O^k}{k!} = 1_{\mathcal{M}_n(\K)} = I_n \quad \text{et} \quad e^{(A+ (-A))} = e^A e^{-A} = I_n \] 
\end{proposition}

\begin{prop}[Dérivation de l'exponentielle matricielle] 
		\begin{itemize} 
			\item Soit $A \in \mathcal{A}_n(\K)$. La fonction:
			\[ 
				\begin{cases} 
					\R \longrightarrow \mathcal{M}_n(\K) \\ 
					t \longmapsto e^{tA} 
				\end{cases}
			\]
			est de classe $ \mathcal{C}^\infty$ sur $\R$ et: 
					\[ \boxed{ \frac{\partial e^{tA}}{\partial t} = A e^{tA} = e^{tA} A } \]
			\item Soit $A \in \mathcal{C}^1(\R, \mathcal{M}_n(\K))$ telle que 
			$ \forall t \in \R, A(t)A'(t) = A'(t) A(t)$  alors: 
				\[ t \longmapsto e^{A(t)} \in \mathcal{C}^1 \quad \text{et} \quad \frac{\partial e^{A(t)}}{\partial t} = A'(t)e^{A(t)} \]
		\end{itemize}
		Le théorème fondamental de l'analyse est toujours vrai. 
\end{prop}

\begin{remark} 
	Soit $f \in \mathcal{C}^1(I,\mathcal{M}_n(\K))$ la formule:
		\[ (f^2)' = 2f'f \] 
	est fausse puisque $( \mathcal{M}_n(\K), +, \times)$ n'est pas un anneau commutatif. 
\end{remark} 

\begin{definition}[Exponentielle] 
	Soit $z \in \C$ on définit l'exponentielle de $z$ par: 
		\[ e^z = \exp (z) = \sum_{k=0}^{\infty} \frac{z^k}{k!} \]
\end{definition} 

\begin{definition}[Exponentielle de Matrice]
	Soit $A \in \mathcal{M}_n(\K)$ on a alors: 
		\[ \exp(A) = e^A = \sum_{k=0}^{\infty} \frac{A^k}{k!} \] 
	Cette série converge toujours. 
\end{definition} 

\begin{remark} 
	Cette définition nous servira plus tard pour la résolution d'équations 
	différentielles en dimension $n \in \N$. 

	On peut remarquer que la fonction $e^A : t \longmapsto e^{A(t)} $ est solution de 
		\[ (E) : Y' = A'Y \]
\end{remark} 

% ==================================================================================================================================
% Théorème de Cauchy-Lipschitz Linéaire 

\section{Théorème de Cauchy Lipschitz Linéaire} 

\begin{definition}[EQ linéaire du premier ordre]
	Soit $A \in \mathcal{C}^0(I, \mathcal{M}_n(\K))$ et $ B \in \mathcal{C}^0(I,\K^n)$. 
	Une équation différentielle linéaire du premier ordre est une équation du type: 
		\[ \boxed{  \forall t \in I, \quad (E) : y'(t) = A(t)y(t) + B(t) = f(t,y)  } \]
		\[ \text{où: } f : 
			\begin{cases} 
				I \times \K^n \longrightarrow \K^n \\ 
				(t,x) \longmapsto A(t)x + B(t) 
			\end{cases} \]
\end{definition} 

\begin{theorem}[Cauchy-Lipschitz Linéaire]
	Soit $(E)$ une équation diférentielle linéaire. Alors ile xiste une unique solution 
	\underline{maximale} au problème de Cauchy: 
		\[ (P) : 
			\begin{cases} 
				(E) \\ 
				y(t_0) = y_0 
			\end{cases} \] 
	avec $(t_0, y_0) \in I \times \K^n$. 
	De plus, cette solution est \textbf{globale}. 
\end{theorem} 


\begin{remark} 
	Pour une équation d'ordre 2 scalaire :
		\[ (E) : y''(t) = a_0(t)y(t) + a_1(t)y'(t) + b(t) \quad \forall t \in U \] 
	où $a_0, a_1,b \in \mathcal{C}^1(I, \K)$ équivalente à $ Y'(t) = A(t) Y(t) + B(t) $ où 
		\[ A(t) =  
			\begin{pmatrix}
				1 & 0 \\ 
				a_0(t) & a_1(t) 
			\end{pmatrix} 
			\quad \text{et} \quad B(t) = 
			\begin{pmatrix}
				0 \\ 
				b(t) 
			\end{pmatrix}
			\quad \text{et} \quad Y(t) = 
			\begin{pmatrix} 
				y(t)  \\
				y'(t) 
			\end{pmatrix} \]
	D'après Cauchy-Lipschitz, il existe une unique solution globale à: 
		\[ 
			\begin{cases} 
				Y'(t) = A(t) Y(t) + B(t) \\ 
				Y(t_0) = Y_0
			\end{cases} 
			\quad \text{où} \quad (t_0, Y_0) \in I \times K^n 
		\]
	Alors il existe une unique solution à l'équation globale d'où: 
		\[ \begin{cases} 
				y''(t) = a_0(t) y(t) + a_1(t) y'(t) + b(t) \\ 
				y(t_0) = y_0 \\ 
				y'(tt_0) = y_1 
			\end{cases} \] 
	Pour l'ordre 2, il faut donc fixer une valeur pour chaque ordre de l'équation. 
	pour l'ordre 3, il faut en fixer 3, etc...
\end{remark}

% ==================================================================================================================================
% Structure des solutions 

\section{Structure de l'ensemble des solutions}

Dans cette section, nous allons essayer de donner une structure à 
l'ensemble des solutions d'une équation différentielle. Cela nous 
permettra d'avoir de meilleures propriétés sur ces solutions et, à partir 
de solutions particulières d'en déduire des solutions plus générales. 

\vspace{0.5cm} 

Soit l'équation différentielle définie sur un intervalle $I$ suivante :
	\[ (L) : y'(t) = A(t) y(t) + B(t) \quad (H) : y'(t) = A(t)y(t) \] 
on appelle $(H)$ l'équation homogène de $(L)$. 
Soient $y_1$ et $y_2$ deux solutions de $(H)$ alors $ \forall \lambda \in \K$, on a: 
	\[ (\lambda y_1 + y_2)'(t) = \lambda A(t)y_1(t) + A(t) y_2(t) = A(t) (\lambda y_1(t) + y_2(t)) \]
Donc c'est aussi une solution de $(H)$. 

\begin{theorem}[Structure]
	Soit $(L)$ une équation différentielle linéaire d'ordre 1 et $(H)$ 
	son équation homogène associée. 
	$S_H$ est un sous-espace vectoriel de $ \mathcal{C}^0(I, \K^n)$ de dimension finie 
	égale à $n$ (on peut donc trouver une base de $S_H$). On a donc: 
		\[ \boxed{ S = S_H + y_1 } \]
	où: 
	\begin{itemize} 
		\item $S$ est l'ensemble des solutions de $(L)$
		\item $y_1$ est une solution particulière de $(L)$
	\end{itemize}
	On dit que $S$ est un espace affine. 
\end{theorem}

\begin{remark}[Construction d'une solution particulière]
		\[ (L) : y(t) = a(t) y(t) + b(t) \quad \forall i \in I \] 
	Posons $ \forall t \in I, y(t) = C e^{\int_{t_0}^{t} a(t) \; ds}$ avec $C \in \K$. 
	Modifions notre potentielle solution telle que: 
		\[ \forall t \in I, \quad y(t) = C(t) e^{\int{t_0}{t} a(s) \; ds} \]
	où $C$ est une fonction dérivable. On a donc que $y$ est dérivable par produit/composition
	et d'après le théorème fondamental de l'analyse. D'où: 
		\begin{align*}
			y'(t) &= C'(t) e^{\int_{t_0}^{t} a(s) \; ds} + C(t) a(t) e^{\int_{t_0}^{t} a(s) \; ds} \\
				  &= a(t)y(t) + b(t) \\
				  & \iff a(t) y(t) + b(t) = C'(t) e^{i(t)} + C(t) a(t) e^{i(t)} \\ 
				  & \iff b(t) = C'(t) e^{i(t)} \\ 
				  & \iff C'(t) = b(t) e^{-i(t)} 
		\end{align*}
	Donc $ \forall t \in I, C(t) = C(t_0) \int_{t_0}^{t} b(s) \; ds $. 
	
	\textbf{Réciproquement:}

	La fonction $ t \longmapsto \left( \lambda \int_{t_0}^{t} b(s) \; ds \right) e^{-i(t)} $ 
	est bien solution de $(L)$. On peut donc en déduire une solution particulière de $(L)$
	en fixant $\lambda$ est $t_0$. De plus: 
		\[ S_H = vect \left( t \longmapsto e^{ \int_{t_0}^{t} a(s) \; ds } \right) 
		= \left\{ t \longmapsto \lambda e^{ \int_{t_0}^{t} a(s) \; ds} \; | \; \lambda \in \K  \right\} \] 
\end{remark}

\begin{theorem}[Résolution Complète]
	Soit $(L) : y'(t) = a(t)y(t) + b(t) $ une équation différentielle linéaire scalaire  d'ordre 1 où 
	$a,b \in \mathcal{C}^0(I, \K^n)$. Alors l'ensemble des solutions de $S$ de $(L)$ est de la forme: 
		\[ S = S_H + \left\{ t \longmapsto \left( \int_{t_0}^{t} b(s) e^{- \int_{t_0}^{t} a(s) \; ds} \right) e^{\int_{t_0}^{t} a(s) \; ds} \right\} \] 
\end{theorem}

% ==================================================================================================================================
% Matrices Fondamentales 

\section{Matrices Fondamentales}

\subsection{Généralités et propriétés} 

Nous savons que l'ensemble des solutions d'une équation différentielle linéaire est un 
espace affine composé d'un sev et d'une solution particulière. 
Le sev $S_H$ est l'ensemble des solutions de l'équation homogène $(H)$ associée à $(L)$ 
on sait qu'il est de dimension $n \in \N$. On peut donc en trouver une base. 

\vspace{0.3cm}

Soit $(H) : y'(t) = A(t) y(t) \quad \forall t \in I$. Une équation homogène d'ordre 1
en dimension $n \in  \N$ et $A \in \mathcal{C}^0(I, \mathcal{M}_n(\K))$.   

\begin{definition}[Système Fondamental]
	Un système fondamental de solutions est une famille finie qui forme une base 
	vectorielle de $S_H$. 
\end{definition} 

\begin{definition}[Matrice Fondamentale]
	Soit $ \mathcal{F} = (y_1, \dots, y_n)$ un système fondamental de $(H)$. 
	Une matrice fondamentale de $(H)$ est une \textbf{fonction} $\Phi$ telle que : 
		\[ \Phi : 
			\begin{cases}
				I \longrightarrow \mathcal{M}_n(\K) \\ 
				t \longmapsto (y_1(t), \dots, y_n(t))
			\end{cases}  \]	
\end{definition} 

\begin{remark}
	Pour tout $ i \in  I$, $\Phi$ est inversible. 
\end{remark} 

\begin{definition}[Wronksien]
	Le wronksien associé à $(H)$ est la fonction : 
		\[ \omega : 
			\begin{cases}
				I \longrightarrow \K \backslash \{0\} \\ 
				t \longmapsto \det (\Phi(t))
			\end{cases} \]
	où $\Phi$ est une matrice fondamentale associée à $(H)$.
	C'est le déterminant de la matrice fondamentale évaluée en $t \in  I$. 
\end{definition} 

\begin{remark} 
	Soit $\Phi$ une matrice fondamentale de $(H)$ et $P \in  GL_n(\K)$ alors $\Phi P$ est aussi 
	une matrice fondamentale de $(H)$. 
\end{remark} 

Soit $\Phi$ une matrice fondamentale de $(H)$ et $y$ une solution particulière de $(H)$.
Alors $ \exists c_1, \dots, n_1 \in  \K$ tels que: 
	\[ \forall i \in  I, \quad y(t) = \sum_{i=1}^{n} c_i y_i(t) \quad \iff \quad y(t) = \Phi(t) 
		\begin{pmatrix}
			c_1 \\ 
			\vdots \\
			c_n 
		\end{pmatrix} \] 

\begin{proposition}
	L'ensemble des solutions de $(H)$ est de la forme: 
		\[ S_H = \{ \Phi C \; | \; C \in  \K^n \} \] 
	où $\Phi$ est une matrice fondamentale de $(H)$. De plus, l'unique solution 
	du problème de Cauchy :
		\[ (P) : \begin{cases}
				(H) \\ 
				y(t_0) = y_0 \quad (t_0, y_0) \in  I \times \K^n
			\end{cases} \]
	est la fonction $ \Phi \Phi^{-1} (t_0) y_0 : t \longmapsto \Phi(t) \Phi^{-1}(t_0) y_0$.  
\end{proposition} 

On peut donc en conclure que dès que l'on connaît une matrice fondamentale 
d'une équation homogène $(H)$, alors on en connaît toutes les solutions. 
Seulement, comment construire une telle matrice ? 

\begin{prop}[Matrice Fondamentale et Base] 
	Soit $\Phi : I \overset{ \mathcal{C}^1}{\longrightarrow} \mathcal{M}_n(\K)$. $\Phi$ est 
	une matrice fondamentale de $(H)$ ssi:
	\begin{enumerate} 
		\item $ \forall t \in  I, \quad \Phi(t) \in  GL_n(\K) $ (famille libre) 
		\item $ \forall t \in  I, \quad \Phi '(t) = A(t) \Phi(t) $ (famille génératrice de $S_H$)
	\end{enumerate} 
\end{prop} 

\subsection{Méthode du polynôme caractéristique} 

Soit $(H_n)$ une équation linéaire homogène d'ordre $n \in  \N$ en dimension 1 telle que : 
	\[ \forall t \in (H_n) : a_n y^{(n)}(t) = a_0 y(t) + a_1 y'(t) + \dots + a_{n-1} y^{(n-1)}(t) \]
où $ a_n, \dots, a_1 \in  \K$. 

\begin{definition}[Polynôme caractéristique]
	Soit $(H_n)$ une équation différentielle homogène d'ordre $ n \in  \N$. 
	On définit le polynôme caractéristique de $(H_N)$ comme le polynôme tel que: 
		\[ P(X) = X^n + a_{n-1} X^{n-1} + \dots + a_1 X + a_0 \] 
	On dira que \textbf{l'équation caractérisitique de} $(H_n)$ est $P(X) = 0$. 
\end{definition} 

\begin{theorem}[Construction des solutions] 
	Plaçons nous dans $\C$. Soient $r_1, \dots, r_p, \; p \leqslant n$ les racines de $P(X)$ de multiplicités 
	respectives $m_1, \dots, m_p$ dans $\C$. Alors: 
		\[ \forall (k,l) \in  \N^2, \; 1 \leqslant k \leqslant p, \; 0 \leqslant l \leqslant m_k -1 
			\quad F_{k,l} : t \longmapsto t^l e^{r_k \times t} \]
	donne une base $(F,k,l)$ de solutions de $(H_n)$. 
\end{theorem} 


% ==================================================================================================================================
% Méthode de variation de la constante 

\section{Méthode de variation de la constante} 

Soient $(L)$ une équation différentielle linéaire d'ordre 1 en dimension $n \in  \N$
et $(H)$ son équation homogène associée telle que: 
	\[ \forall t \in I, \quad (L) : y'(t) = A(t) y(t) + B(t) \quad y'(t) = A(t) y(t) \] 

Pour rappel, l'ensemble des solutions de cette équation est de la forme: 
	\[ S = S_H + y \] 
où $y$ est une solution particulière de $(L)$ et $S_H$ l'ensemble des solutions de $(H)$. 

Ici, on cherche simplement \textbf{une solution particulière} $y$ de la forme: 
	\[ y = \Phi C \quad \text{où} \; C : \overset{ \mathcal{C}^0}{\longrightarrow} \K^n \] 
et $\Phi$ est la matrice fondamentale associée à $(H)$. 

$y$ est solution donc nécessairement, $ \forall t \in I$: 

\begin{align*}
	y'(t) &= A(t) \Phi(t) C(t) + \Phi(t) C'(t) \\ 
			&= A(t) y(t) + B(t) \\ 
			&= A(t) \left( \Phi(t) C(t) \right) + B(t)  \\ 
\end{align*}

\[ \iff 
\begin{cases}
	\Phi'(t) = A(t) \Phi(t) \\ 
	C'(t) = \Phi^{-1}(t) B(t)
\end{cases} \] 

D'où après intégration : $ C(t) = C(t_0) + \int_{t_0}^{t} \Phi^{-1}(s) B(s) \; ds $. 
Une solution particulière de $(L)$ est donc :

	\[ \boxed{ y : t \longmapsto \Phi(t) \left( C(t_0) + \int_{t_0}^{t} \Phi^{-1}(s) B(s) \; ds \right) } \] 






